A search over every pair of per-group thresholds (an exhaustive walk over the two groups' ROC staircases for the oracle curve, and a grid over the share of the budget given to \texttt{.vn} for the calibrated curve), scored against the budget identity $\mathrm{FPR} = \pi_{vn}\mathrm{FPR}_{vn} + \pi_{other}\mathrm{FPR}_{other}$ with $\pi_{vn} = 0.612$, leaves the conclusion intact and makes it exhaustive rather than illustrative (Table~\ref{tab:budgetfrontier}, Figure~\ref{fig:budgetfrontier}). At the deployed false-alarm budget of $0.259$ the best allocation that holds \texttt{.vn} misses to $10\%$ concedes $0.402$ of non-\texttt{.vn} phishing, and the most balanced point the budget admits still misses $0.284$ of \emph{both} groups; bringing both below $10\%$ requires a corpus false-alarm rate of $0.421\,\pm\,0.006$, $\times1.6$ the deployed budget. These frontiers are \emph{oracle} bounds (their thresholds are read off the test curves, which no deployment can do), so the impossibility they state is the generous form of it: what a calibrated rule reaches (the second column of Table~\ref{tab:budgetfrontier}) is weaker still, and the published rules of Table~\ref{tab:groupthr} sit on that reachable curve rather than inside it. The blind spot is thus not one badly-chosen operating point among good ones: at this budget there is no good one%
